\section{Critical clusters remain controlled after deletions}
\label{good:section}

The local maps constructed in the next section attach pieces cut from
critical clusters. We need a property of the original cluster that
controls the separating bridges of every such piece, even after an
arbitrary vertex deletion. This section proves that a capped critical
cluster fails this property only with stretched-exponentially small
probability. All percolation fields in this section have parameter
$\pc$.

\begin{definition}[Pivotal distance and good clusters]
\label{good:definition}
A bridge of a connected open graph is an open edge whose deletion
disconnects that graph. The \emph{pivotal distance} between two vertices
is the number of bridges separating them. Equivalently, delete all
bridges, contract the resulting connected pieces, and measure distance
in the resulting tree. Write $\operatorname{diam}_{\mathrm{piv}}$ for
the maximum pivotal distance and
$\operatorname{rad}_{\mathrm{piv}}(H,x)$ for its maximum from $x$.

Fix $0<\epsilon<1/2$ and set $s=1/2+\epsilon$. A finite open cluster
$C$ with $|C|\le M$ is \emph{good for cap $M$} if the following holds
for every set $A\subseteq C$. For every open component $B$ after
deleting $A$, let $K$ be the number of original open edges from $B$
to $A$. Whenever $K\ge1$,
\[
 \operatorname{diam}_{\mathrm{piv}}(B)\le K M^s.
\]
The pivotal distance in this condition is computed inside $B$.
\end{definition}

\begin{proposition}[Good critical clusters]\label{good:good-clusters}
For each $\epsilon\in(0,1/2)$, uniformly in $x$ and for all sufficiently
large $M$,
\begin{equation}\label{eq:s06}
 \Prob_\pc(s_x\le M,\ C_x\text{ is not good for cap }M)
 \le e^{-M^\epsilon}.
\end{equation}
The event that a cluster has size at most $M$ and is good for cap $M$
is determined by the percolation bits in a deterministic finite-radius
ball about its root.
\end{proposition}

We prove the proposition through a marked-cluster estimate. Independent
vertex marks, often called a ghost field, are a standard device in
percolation differential inequalities \cite{AizenmanBarsky1987}.
The exploration-score method is closely related to that used for the
two-ghost inequality in \cite[Section~3]{Hutchcroft2020Locality};
the deletion-uniform estimate below requires the additional inverse
comparison at the end of this section. In a
deterministic induced domain $D$ containing $x$, mark vertices
independently with probability $h\in(0,1/2]$, independently of the
edges. Let $E=E(D,x,h)$ be the event that $C_x^D$ contains a mark.
An edge is pivotal for $E$ if changing its state changes $\one_E$,
with the marks held fixed. Let $D_p$ denote the number of open
pivotal edges; this notation is confined to the present section and
the percolation parameter is $p=\pc$.

\begin{lemma}[Marked pivotal moments]\label{good:marked-pivotals}
There is a constant $C$, depending only on $d$ and $\pc$, such that,
with
$A_h=C\sqrt{\log(1/h)/h}\ge1$, for every induced domain $D$, every
$x\in D$, every integer $i\ge0$, and every $r\ge0$,
\begin{equation}\label{eq:s07}
 \begin{aligned}
 \E\left[\binom{D_p}{i}\one_E\right]
 &\le A_h^i\Prob(E),\\
 \Prob(E,\ D_p\ge r)
 &\le2\Prob(E)e^{-r/(4A_h)}.
 \end{aligned}
\end{equation}
\end{lemma}
\begin{proof}
\emph{The first moment.}
Explore the cluster in a fixed mark-independent order, testing the
mark of each newly discovered vertex immediately and stopping at the
first mark or when the cluster is exhausted. No edge is queried twice.
Critical finiteness from \cref{pre:critical-finiteness} makes this a
terminating decision procedure, even in an infinite domain.
Let $T$ be its number of edge queries. Conditional on the complete
percolation configuration, the discovery order is fixed. The first
mark index is geometric with parameter $h$, stopped at $|C_x^D|$;
conditional on success, it is that geometric variable conditioned to
be at most $|C_x^D|$. In both cases its mean is at most $1/h$.
There are at most $d$ edge queries per discovered vertex, so
\[
 \E T\le d/h,\qquad \E[T\mid E]\le d/h.
\]

Let $Z=\sum_{e\text{ queried}}(\omega_e-\pc)$. The indicator $Q_e$
that edge $e$ is queried is measurable without its own bit, and every
pivotal edge must be queried by a procedure deciding $E$. Conditioning
on all other edge bits and marks, and then averaging, gives
\[
 \E[(\omega_e-\pc)Q_e\one_E]
 =\pc(1-\pc)\Prob(e\text{ is pivotal}).
\]
The identity remains valid after summation, since $\E T<\infty$.
An open pivotal contributes to $D_p\one_E$, and its bit is independent
of its pivotality. Therefore
\begin{equation}\label{good:score-identity}
 \E[Z\one_E]=(1-\pc)\E[D_p\one_E].
\end{equation}

Each queried bit is fresh conditional on the previous transcript.
For a centered Bernoulli variable, the logarithmic moment-generating
function is at most $\lambda^2/8$: its second derivative is a
Bernoulli variance, at most $1/4$. Successive conditional
expectations at the bounded stopping times $T\wedge n$, followed by
Fatou's lemma, give
\[
 \E\exp(\lambda Z-\lambda^2T/8)\le1.
\]
Put $H=\Prob(E)\ge h$, since a mark at $x$ suffices. Conditional
Jensen's inequality yields
\[
 \lambda\E[Z\mid E]
 \le\log(1/H)+\frac{\lambda^2}{8}\E[T\mid E]
 \le\log(1/h)+\frac{\lambda^2d}{8h}.
\]
Taking $\lambda=\sqrt{8h\log(1/h)/d}$ and using
\eqref{good:score-identity}, we obtain
\[
 \E[D_p\mid E]
 \le\frac{\sqrt{d/2}}{1-\pc}
       \sqrt{\frac{\log(1/h)}h}.
\]
This proves the first-moment assertion, uniformly over domains.

\emph{Higher binomial moments.}
On $E$, each open pivotal is a bridge separating $x$ from all marks
in its cluster. These pivotals are linearly ordered from $x$ toward
the marks. Define
\[
 b_i(D,x)=\E\left[\binom{D_p}{i}\one_E\right],
 \qquad b_0(D,x)=\Prob(E).
\]
For an oriented edge $\vec e=(a,b)$, let $C_{x,\vec e}^D$ be the
root cluster when $e$ is forced closed. Decompose a selected set of
$i$ pivotals according to its first selected edge, oriented away
from $x$. The resulting exact identity is
\begin{equation}\label{good:pivotal-recursion}
 b_i(D,x)=\pc\sum_{\vec e=(a,b)\in E(D)}
 \sum_{\substack{W\subset D\text{ finite}\\x,a\in W, b\notin W}}
 \Prob_\pc(C_{x,\vec e}^D=W)(1-h)^{|W|}
 b_{i-1}(D\setminus W,b).
\end{equation}
To verify the conditional law in this formula, force $e$ closed and
expose the exact finite root cluster $W$. This reveals internal and
boundary information of $W$ but no edges internal to $D\setminus W$.
The event $C_{x,\vec e}^D=W$ ignores the actual bit of $e$.
Require that $W$ contain no marks, accounting for $(1-h)^{|W|}$.
Opening $e$, with independent probability $\pc$, makes it the sole
open contact from $W$ to its exterior. It is then pivotal exactly
when $b$ connects to a mark in $D\setminus W$. All pivotals after
it are precisely those of that exterior event. The exterior edges
and marks retain their independent laws. Every selected subset is
counted once: if the total number of ordered pivotals is $n$, the
count is $\sum_{j=1}^n\binom{n-j}{i-1}=\binom ni$.
All sums are nonnegative, so Tonelli's theorem justifies the
decomposition without a priori moment assumptions.

Uniform induction in $D,x$ now gives
\[
 b_i(D,x)\le A_h^{i-1}b_1(D,x)\le A_h^ib_0(D,x).
\]
Finally expand $(1+(2A_h)^{-1})^{D_p}$ and sum these binomial
moments to get
\[
 \E\left[\left(1+\frac1{2A_h}\right)^{D_p}\one_E\right]
 \le2\Prob(E).
\]
Since $\log(1+(2A_h)^{-1})\ge(4A_h)^{-1}$ for $A_h\ge1$,
Markov's inequality proves the second part of \eqref{eq:s07}.
\end{proof}

\begin{lemma}[A capped pivotal-radius tail]\label{good:radius-tail}
For the fixed $\epsilon\in(0,1/2)$, there is $c>0$ such that, uniformly
in $x$, all sufficiently large $M$, and integers
$1\le K\le M^{1-s}$,
\begin{equation}\label{eq:s08}
 \Prob_\pc\left(s_x\le M,
     \operatorname{rad}_{\mathrm{piv}}(C_x,x)\ge KM^s/2\right)
 \le\exp\left(-\frac{cK^2M^{2\epsilon}}{\log M}\right).
\end{equation}
\end{lemma}
\begin{proof}
On the event on the left, choose a witness vertex $y$ by a fixed
ordering. Independently mark vertices with probability $h$. The
conditional probability that $y$ is the sole mark in $C_x$ is at
least $h(1-h)^M$. In that case every bridge separating $x$ from $y$
is pivotal for the marked event. By \eqref{eq:s07}, with $R=KM^s/2$,
\[
 \Prob_\pc(s_x\le M,\operatorname{rad}_{\mathrm{piv}}\ge R)
 \le\frac{2}{h(1-h)^M}\exp\left(-\frac{R}{4A_h}\right).
\]
Set
\[
 X=\frac{K^2M^{2\epsilon}}{\log M},\qquad
 h=\frac{X}{bM}
   =\frac{K^2M^{-1+2\epsilon}}{b\log M},
\]
where $b$ will be a sufficiently large fixed constant. The allowed
range of $K$ gives $h\le1/(b\log M)\le1/2$. For all sufficiently
large $M$, uniformly in that range,
$\log(1/h)\le2\log M$ and
\[
 \frac{R}{4A_h}\ge\frac{X}{8C\sqrt{2b}}.
\]
Using $-\log(1-h)\le2h$, the logarithm of the preceding probability
bound is therefore at most
\[
 O(\log M)+\frac{2X}{b}-\frac{X}{8C\sqrt{2b}}.
\]
First choose $b$ large enough that the last two terms have a negative
sum bounded above by $-c_1X$. Then increase $M$: the minimum possible
$X$ is $M^{2\epsilon}/\log M$, which dominates $\log M$. This proves
\eqref{eq:s08} with a fixed positive $c$.
\end{proof}

\begin{proof}[Proof of \cref{good:good-clusters}]
Given a capped cluster violating goodness, choose a deleted set $A$
and a remaining component $B$ witnessing the violation, by a fixed
finite ordering. Let $K\ge1$ be the number of original open edges
from $B$ to $A$. Since
$\operatorname{diam}_{\mathrm{piv}}(B)\le|B|-1<M$, a violation has
$K\le M^{1-s}$.

Close its $K$ open contacts, except for one retained contact if
$x\notin B$: in that case take a simple open path from $x$ to $B$
and retain its first edge entering $B$. The path before that contact
does not use any edge being closed, so the modified root cluster
$C'_x$ still contains $B$. If $x\in B$, then $C'_x=B$. In either
case $B$ has at most one open attachment to the rest of $C'_x$.
An internal bridge of $B$ cannot be bypassed outside $B$, because
such a bypass would need two attachments. Thus its bridge tree
embeds into the bridge tree of $C'_x$ with its distances preserved.
The triangle inequality shows that at least one endpoint of a
diameter of $B$ is at pivotal distance greater than $KM^s/2$ from
$x$. Also $|C'_x|\le|C_x|\le M$. Hence the modified configuration
belongs to the event in \eqref{eq:s08}.

It remains to bound the cost of this configuration change; no
independence of its input-dependent choices is needed. At a fixed
output, every changed edge is incident to $B\subseteq C'_x$, so
there are at most $dM$ candidate edges. For fixed $K$, specifying
the subset of at most $K$ changed edges reconstructs the input by
reopening them. The number of possible inputs is at most
\[
 \sum_{j=0}^K\binom{\lfloor dM\rfloor}{j}\le(K+1)(dM)^K,
\]
and their probability ratio to the output is at most
$\max(1,\pc/(1-\pc))^K$. This is a comparison on a finite product:
on the cap event every cluster vertex lies within distance $M-1$
of $x$, and every incident edge lies in a deterministic slightly
larger ball. Condition on all bits outside that ball and apply the
pointwise finite-product comparison inside it.

It follows from \eqref{eq:s08} that the contribution for fixed $K$
is at most
\[
 \exp\left(C_1K\log M-
                 \frac{cK^2M^{2\epsilon}}{\log M}\right).
\]
Uniformly for $K\ge1$, the negative term absorbs the positive one,
because $M^{2\epsilon}/(\log M)^2\to\infty$. Summing over the at
most $M$ possible values of $K$ gives, for large $M$,
\[
 M\exp\left(-\frac{cM^{2\epsilon}}{2\log M}\right)
 \le e^{-M^\epsilon},
\]
as claimed. Finally, explore the root cluster until it is exhausted
or more than $M$ vertices are found. This uses only a deterministic
ball of radius at most $\lceil M\rceil+1$. In the first case its
internal edges and closed boundary are known, and one can test
every subset deletion of the finite cluster. In the second case
the capped event fails. This proves locality.
\end{proof}
